\section{The Reporting Card}\label{sec:card}\label{sec:methodology}
The three rules become usable only if a reader can apply them to a specific paper and a specific deployment. The reporting card (Table~\ref{tab:r33-card}) records, for one source configuration, the evidence for each boundary together with the function and cost measured under it.

\begin{table}[t]\centering\footnotesize
\caption{The reporting card. Fill one card per source configuration and proposed deployment. Every evidence cell carries a source version and section/table locator and one status code.}\label{tab:r33-card}
\begin{tabularx}{\columnwidth}{@{}p{.21\columnwidth}Y@{}}\toprule
Field & What to enter\\\midrule
Receiver/target & Hidden property; the signal and fields the observer receives; observer permissions.\\
Intervention & Mechanism; where it executes (stored trace, wire, tool action); whether cover is discarded, isolated or executed.\\
Attacker/test & Data the attacker or calibration was fitted on; defended test population; held-out unit.\\
Function & Preservation witness, output quality or task outcome, at the configuration used for protection.\\
Cost & Resource, unit, baseline, denominator and included operations.\\
Extension & What the deployment changes (model, workload, hardware, network); which of these the source held out.\\\midrule
\multicolumn{2}{@{}p{.96\columnwidth}@{}}{\textbf{Status codes.} \textbf{B} bound: every element the claim needs is stated for this configuration. \textbf{P} partial: some are stated, at least one is not. \textbf{U} not located in the inspected sections. \textbf{?} stated but not verified. \textbf{--} not applicable to the claim.}\\
\bottomrule\end{tabularx}\end{table}


\subsection{From fields to a decision}
A reuse decision starts from a \emph{claim}, not from a paper: ``this defense protects my deployment's TTFT channel'', ``its quality loss is negligible on my model'', ``its overhead is 6\,ms''. Each claim depends on a subset of the fields. A protection claim needs Receiver, Intervention, Attacker and Extension; a cost claim needs Intervention, Cost and Extension; a function claim for executed cover needs Intervention and Function at the task level. The decision rule is:
\begin{itemize}
\item \textbf{Supported} if every required field is \textbf{B} and the deployment matches it.
\item \textbf{Incompatible} if a required field shows a demonstrated mismatch, for example a cost measured with the protection disabled.
\item \textbf{Unresolved} otherwise. The output names the fields that are \textbf{P}, \textbf{U} or \textbf{?}, and therefore the measurement still needed.
\end{itemize}
Figure~\ref{fig:main-decision} shows the procedure. Three outcomes may seem coarse, but the information lies in their combination. One paper rarely yields a single verdict; it yields different outcomes for different claims, and an unresolved outcome is actionable because it specifies what to measure.

\begin{figure*}[t]\centering
\begin{tikzpicture}[font=\scriptsize,>=Latex,
check/.style={draw,rounded corners,fill=blue!6,align=left,text width=4.8cm,minimum height=16mm},
test/.style={draw,fill=orange!10,align=left,text width=4.8cm,minimum height=16mm}]
\node[check] (receiver) at (0,0) {\textbf{1. Bind the receiver}\\Target, actual fields and linkage; does protection reach the final emitter?};
\node[check] (action) at (5.65,0) {\textbf{2. Bind the execution boundary}\\Is cover discarded, isolated, or admitted to tool/GUI execution?};
\node[check] (effect) at (11.3,0) {\textbf{3. Bind the evaluation}\\Attacker access/adaptation, workload, inference, utility and cost object.};
\draw[->] (receiver)--(action);
\draw[->] (action)--(effect);
\node[test] (rtest) at (0,-2.0) {Changed emitter or view:\\capture the final hop.\\Same verified projection: retain the source-local inference result.};
\node[test] (atest) at (5.65,-2.0) {Executed cover: check terminal goals and policy. Verified discard excludes direct action effects; retain delay/resource checks.};
\node[test] (etest) at (11.3,-2.0) {Changed mechanism or workload:\\measure that component/condition.\\Adaptation claims require the same defended-test comparison.};
\draw[->] (receiver)--(rtest);
\draw[->] (action)--(atest);
\draw[->] (effect)--(etest);
\node[draw,rounded corners,fill=green!8,align=center,text width=16.1cm,minimum height=10mm] (out) at (5.65,-3.7)
{\textbf{Output for the requested claim:} supported within matched conditions; incompatible because of an observed mismatch; or unresolved because a required binding is missing. No branch certifies universal protection.};
\draw[->] (rtest.south)--++(0,-.25)-|(out.north west);
\draw[->] (atest)--(out);
\draw[->] (etest.south)--++(0,-.25)-|(out.north east);
\end{tikzpicture}

\caption{Decision procedure for reusing a defense result. Each step binds one boundary; the output is per claim.}\label{fig:main-decision}
\end{figure*}

\subsection{A complete example}
Table~\ref{tab:worked-card} fills every field for CachePrune \cite{wu2026cacheprune}, which excludes detector-marked sensitive spans from cross-user KV-cache sharing. The hypothetical deployment runs a different model and its tenants can time responses remotely. The card separates three conclusions that a summary such as ``CachePrune protects prompts at little cost'' would merge. The recovery results support protection against binary hit/miss leakage, because the source's abstract oracle is at least as strong as a remote timer for that signal. The fixed-chunk comparison in \S5.7 cannot price protection, because it disables privacy filtering. Quality, latency and detector recall must be measured again on the target. Filling the card required CachePrune's \S3.1, \S5.1--5.7 and Tables~1--2.

\begin{table*}[t]\centering\footnotesize
\caption{A complete card: reusing CachePrune v1 \cite{wu2026cacheprune} for a hypothetical multi-tenant service that runs a different model and whose tenants measure time to first token (TTFT) remotely. One source yields all three outcomes, each tied to a specific claim.}\label{tab:worked-card}
\setlength{\tabcolsep}{3pt}
\begin{tabularx}{\textwidth}{@{}>{\raggedright\arraybackslash}p{.09\textwidth}Yp{.045\textwidth}>{\raggedright\arraybackslash}p{.16\textwidth}Y@{}}\toprule
Field & Source evidence (locator) & Code & Proposed deployment & Outcome for this claim\\\midrule
Receiver/target & Abstract binary cache-reuse oracle; exact recovery of detector-marked sensitive spans on QASPER (\S3.1, Table~1). & B & Remote tenants time TTFT over a WAN. & \textbf{Supported} for binary hit/miss leakage: a perfect oracle reveals at least what a noisy timer does. \textbf{Unresolved} for timing that reveals more than hit/miss, such as partial-hit length.\\
Intervention & Detector-marked spans are excluded from cross-user KV-cache sharing (\S3.1). Detector recall on other data not reported. & P & Same mechanism, different user data. & \textbf{Unresolved}: protection is conditional on the detector; measure its recall on deployment data.\\
Attacker/test & Direct recovery 0\%, contextual recovery 3.21\% (Table~1). No learned attacker, so a clean-versus-defended fit is inapplicable; broader adaptive attacks not located. & B / -- / U & Same two attacks; adaptive tenants possible. & \textbf{Supported} for the two evaluated attacks; \textbf{unresolved} for attacks that exploit unmarked context.\\
Function & QASPER F1 0.28 vs.\ 0.30 without sharing, Mistral-7B, default configuration (Table~2). Same configuration as Table~1, but paired runs not established. & P & Different model. & \textbf{Unresolved}: re-measure quality on the target model.\\
Cost & TTFT 273 vs.\ 797\,ms without sharing, one H100 (Table~2). The fixed-chunk comparison (\S5.7) disables privacy filtering. & P & Different model and hardware. & \S5.7: \textbf{incompatible} as a cost of the protected configuration. Table~2: \textbf{unresolved} for the target hardware.\\
Extension & One model, one GPU type, one dataset; no model or hardware holdout located. & U & Model, hardware and network change. & Measure on the target.\\\midrule
\multicolumn{5}{@{}p{.97\textwidth}@{}}{\textbf{Decision.} Retain the oracle-level recovery result as protection evidence for binary hit/miss leakage. Do not cite \S5.7 as the cost of protection. Before deployment, measure detector recall on deployment data, quality and TTFT on the target model and hardware, and leakage from timing features beyond hit/miss. Inputs: CachePrune v1 \S3.1, \S5.1--5.7, Tables~1--2.}\\
\bottomrule\end{tabularx}\end{table*}


\subsection{Applying the card to defense papers}
Table~\ref{tab:r42-defense-first} applies the card to six defense papers. All five for which output quality applies report a functional score, and the sixth (CacheSolidarity) supports exact-prefix preservation, so the card's role is not to find missing measurements. It shows which measurements may be joined. OCELOT binds protection, task success and an 18.3\% wall-clock overhead at the same release boundary \cite{xie2026ocelot}, so the three can be reused together. PlanTwin's re-identification diagnostic uses unrestricted fields, so it cannot be joined to planning quality under the privacy budget \cite{yu2026plantwin}. DP-MemView's AUC of 0.976 applies to the exposed-score view only (Rule~1). SICC and OCELOT evaluate adaptive attacks beyond the clean-versus-defended fitting contrast; recording them under Attacker/test keeps that evidence visible rather than reducing it to a single adaptive/non-adaptive flag.

\begin{table*}[t]\centering\footnotesize
\caption{The card applied to six defense papers. Each row records which receiver the protection result covers, which function was measured with it, and what limits the reuse of its cost or transfer evidence. Versions and section locators are in the supplement (Post-hoc Checks).}\label{tab:r42-defense-first}
\setlength{\tabcolsep}{3pt}
\begin{tabularx}{\textwidth}{@{}p{.14\textwidth}p{.31\textwidth}p{.20\textwidth}Y@{}}\toprule
Source & Receiver / protection evidence & Function measured & Cost / transfer boundary\\\midrule
CacheSolidarity~\cite{pennas2026cachesolidarity} & Cross-user TTFT; 20-probe timing demonstration, not a scored recovery rate. & Exact-prefix preservation (logical). & TTFT, hit rate and component costs; model and load varied, calibration not held fixed.\\
PlanTwin~\cite{yu2026plantwin} & Cloud-visible planning schema; re-identification measured with unrestricted fields, not under the privacy budget. & Planning quality. & Timing of heuristic default and small-model variant needs a configuration match.\\
SICC~\cite{xu2026sicc} & Form-only (length/structure) inference isolated; adaptive transcript attacks. & Protocol completion per receiver. & No compiler LLM calls; runtime cost partly measured.\\
DP-MemView~\cite{dpmemview2026} & Semantic transcript inference at the compliant interface; separate exposed-score diagnostic (AUC 0.976). & Response quality and personalization. & Runtime cost not bound; corpus transfer fixes scaffold and auditor.\\
OCELOT~\cite{xie2026ocelot} & Semantic trajectory inference; held-out stress proxies and adaptive attacks. & Task success, same condition. & 18.3\% wall-clock overhead at the same boundary.\\
AirGapAgent~\cite{bagdasarian2024airgap} & Contextual extraction by a third party, not passive metadata inference. & Appropriate-answer utility. & Runtime cost not bound; model and scenario variation, calibration not held fixed.\\
\bottomrule\end{tabularx}
\end{table*}


\subsection{Filling, disagreement and adoption}
\paragraph{Procedure}
(1) State the claim and the deployment change. (2) For each required field, quote the source with a locator, or write ``not located''. (3) Assign a status code. (4) Derive the outcome and, if it is unresolved, the missing measurement. Inspected sections are recorded, so a \textbf{U} means ``not located in these sections'', which another reader can check.

\paragraph{Disagreement}
When two readers fill the same card, we propose that each cell take the weaker code (\textbf{B} $>$ \textbf{P} $>$ \textbf{U}) until one reader produces a locator that settles it; a disagreement that remains is reported in the card rather than resolved by authority. This rule is a proposal; its effect on agreement has not been measured (Section~\ref{sec:validation}).

\paragraph{Who uses the card}
Defense authors can include a filled card as an appendix, stating which receiver, execution mode and attacker fit each number belongs to. Reviewers can ask three questions---which emitter's trace was measured, does cover execute, and was the attacker fitted on defended outputs---that correspond to Rules~1--3. Benchmark designers can log the same fields with every result, so that later reuse does not depend on reconstructing them from prose. An editable template accompanies the artifact.
